% A0 portrait poster. Change a0paper to a1paper/a2paper, or landscape for a wide one.
\documentclass[30pt,a0paper,portrait]{tikzposter}
\usepackage[T1]{fontenc}
\usepackage{lmodern}
\usepackage{amsmath,amssymb}
\usepackage{booktabs}
\usepackage{graphicx}

\usetheme{Default}
\usecolorstyle[colorOne=white,colorTwo=blue!10,colorThree=blue!30]{Default}
\usebackgroundstyle{Default}
\colorlet{blocktitlefgcolor}{black}
\useblockstyle{Default}

\title{Poster Title: One Clear Message in Two Lines}
\author{First Author$^1$, Second Author$^2$}
\institute{$^1$Example University \quad $^2$Example Institute\\
  \texttt{first.author@example.edu}}

\begin{document}
\maketitle

\begin{columns}
  \column{0.5}
  \block{Motivation}{
    Say why anyone should stop at your poster, in two or three short lines:
    \begin{itemize}
      \item The problem in one sentence.
      \item Why existing approaches fall short.
      \item What you did instead.
    \end{itemize}
  }

  \block{Method}{
    Use a diagram rather than a paragraph. A formula should be big enough to read
    from two metres away:
    \[ \hat{y} = \operatorname*{arg\,max}_{c}\; p(c \mid x) \]
  }

  \block{Data}{
    Describe the data or the experiment in two or three lines, and give the one number
    that matters: how many samples, which instruments, which period.
  }

  \column{0.5}
  \block{Results}{
    \centering
    \begin{tabular}{lcc}
      \toprule
      Method & Accuracy & Speed-up \\
      \midrule
      Baseline & 81\% & 1$\times$ \\
      Ours & \textbf{87\%} & \textbf{3.2$\times$} \\
      \bottomrule
    \end{tabular}
    \\[1cm]
    \tikzset{every picture/.style={}}
    % \includegraphics[width=0.9\linewidth]{figure.pdf}
    Add a figure here: a large plot with a one-line caption beats a table.
  }

  \block{Conclusions}{
    \begin{itemize}
      \item The main finding, as a full sentence.
      \item The second finding.
      \item What comes next.
    \end{itemize}
  }
  \block{Next steps}{
    \begin{itemize}
      \item What you would do with more time or data.
      \item Where to find the code, the data and the full paper: put a QR code or URL here.
    \end{itemize}
  }
\end{columns}

\block[titlewidthscale=0.5]{References}{
  \small
  [1] L. Lamport, \emph{\LaTeX: A Document Preparation System}, Addison-Wesley, 1994.\\[0.3em]
  [2] D. E. Knuth, ``Literate programming'', \emph{The Computer Journal} 27(2), 1984.
}

\end{document}
